\input pdfToolbox
\input pdfmsym
\pdfmsymsetscalefactor{10}
\input ../preamble
\def\enumindent{0cm}

\def\printmcount{\the\counter{section}.\the\counter{math counter}}
\setcounter{section}{11}

\centerline{\setfontandscale{bf}{20pt}Commutative Algebra \& Algebraic Geometry}
\medskip
\centerline{\bf Exercise \the\counter{section}}
\smallskip
\centerline{\it Ari Feiglin}

\bprob

    \benum
        \item Let $f_i\colon X_i\to Y_i$ be morphisms of varieties for $i=1,2$.
        \benum
            \item Show that there is a canonical morphism $f_1\times f_2\colon X_1\times X_2\to Y_1\times Y_2$.
            \item Show that if $f_1,f_2$ both satisfy any of the below properties, so too does $f_1\times f_2$:

            {\tabskip=0pt plus 1fil\halign to\hsize{(#)\hfil\tabskip=.25cm&#\hfil\tabskip=0pt plus.2fil&(#)\hfil\tabskip=.25cm&#\hfil\tabskip=0pt plus 1fil\cr
                1&being an isomorphism,&2&being an open embedding,\cr
                3&being a closed embedding,&4&being affine.\cr
            }}
        \eenum
        \item Show that $\bA^n\times\bA^m\cong\bA^{n+m}$.
        \item Let $X$ and $Y$ be nonempty algebraic varieties.
        \benum
            \item Show that $\dim(X\times Y)=\dim X+\dim Y$.
            \item Prove that $X$ and $Y$ are connected iff $X\times Y$ is.
            \item Prove that if $X\times Y$ is irreducible, then $X$ and $Y$ are.
            \item Assume that $X$ and $Y$ are irreducible.
            \benum
                \item Suppose that $X\times Y=C_1\cup C_2$ where $C_1,C_2$ are proper closed subsets of $X\times Y$.
                Show that for any $y\in Y$, either $X\times\set y\subseteq C_1$ or $C_2$.
                \item For $i=1,2$, define $F_i=\set{y\in Y}[X\times\set y\subseteq C_i]$.
                Show that $F_i=\bigcap_{x\in X}\pi_Y(C_i\cap(\set x\times Y))$ and that $F_i$ is closed.
                \item Conclude that $X\times Y$ is irreducible.
            \eenum
            \item Prove that $X$ and $Y$ are separated iff $X\times Y$ is.
            \item Prove that $X$ and $Y$ are complete iff $X\times Y$ is.
        \eenum
    \eenum

\eprob

\benum
    \item
    \benum
        \item In any category with all limits of a certain shape, taking limits of that shape is functorial.
        In this case, taking finite products is functorial.
        A natural transformation between product diagrams amounts to morphisms between the objects whose product we are taking, without restrictions.

        Because the set underlying the product of varieties is the product of the underlying sets of varieties, the underlying set-theoretic function of
        $f_1\times f_2$ is $(x,y)\mapsto(f_1(x),f_2(y))$.

        \item For (1) this is due to functorality.
        We see that $(f_1\times f_2)\circ(g_1\times g_2)=(f_1\circ g_1)\times(f_2\circ g_2)$.
        Thus
        $$ (f_1\times f_2) \circ (f_1^{-1}\times f_2^{-1}) = \id\times\id = \id , $$
        and we have a similar equality in the other direction, so $f_1\times f_2$ is an isomorphism.

        For (2) and (3), we claim that if $f_1,f_2$ are embeddings then so is $f_1\times f_2$.
        We then see that restricting $f_1\times f_2$ to $X_1\times X_2\to f_1(X_1)\times f_2(X_2)$ is an isomorphism from (1),
        and since the induced SWF structure of the product and the product of induced structures is the same, this means that $f_1\times f_2$
        is indeed an embedding.

        In lecture we showed that the product of open subsets of a variety is open, and the product of closed subsets is closed.
        Thus in the case of (2), $f_1(X_1)\times f_2(X_2)$ is open, and it is closed in (3).
        Thus $f_1\times f_2$ is an open embedding in (2) and closed in (3).

        For (4), if $U_i$ is an open affine cover of $Y_1$ and $V_j$ an open affine cover of $Y_2$, then $U_i\times V_j$ is an open affine cover of $Y_1\times Y_2$.
        This is because the product of affine varieties is affine, and the product of open subsets if open.
        For each $i,j$, we have that
        $$ (f_1\times f_2)^{-1}(U_i\times V_j) = (f_1^{-1}U_i)\times(f_2^{-1}V_j) , $$
        and since $f_1,f_2$ are affine, $f_1^{-1}U_i$ and $f_2^{-1}V_j$ are affine.
        Thus this preimage is affine, as desired.
    \eenum

    \item Recall that for an SWF $X$,
    $$ \hom(X,\bA^{n+m}) \cong \O_X(X)^{n+m} \cong \O_X(X)^n \times \O_X(X)^m \cong \hom(X,\bA^n) \times \hom(X,\bA^m) , $$
    which is the universal property of the product $\bA^n\times\bA^m$.
    By Yoneda, $\bA^{n+m}\cong\bA^n\times\bA^m$.

    \item
    \benum
        \item If $X,Y$ are affine of dimensions $n=\dim X$ and $m=\dim Y$, then there are finite surjections $f\colon X\to\bA^n$ and $g\colon Y\to\bA^m$.
        The product map $X\times Y\to\bA^{n+m}$ is then surjective, and we claim finite.
        Indeed, $k[X\times Y]\cong k[X]\otimes_kk[Y]$, and if $k[X]$ is a finite $k[\bA^n]$-algebra and $k[Y]$ a finite $k[\bA^m]$-algebra, then
        $k[X]\otimes k[Y]$ is a finite $k[\bA^n]\otimes k[\bA^m]\cong k[\bA^{n+m}]$-algebra.

        Thus $\dim(X\times Y)=n+m=\dim X+\dim Y$.

        Now in general, let $U_i$ be an open affine cover of $X$ and $V_j$ an open affine cover of $Y$.
        Then $U_i\times V_j$ is an open affine cover of $X\times Y$, so
        $$ \multline{
            \dim(X\times Y) = \sup_{i,j}\dim(U_i\times V_j) = \sup_{i,j}(\dim U_i+\dim V_j) = (\sup_i\dim U_i) + (\sup_j\dim V_j)\cr
            = \dim X + \dim Y ,
        } $$
        where the second equality is due to what we just proved.

        \item Suppose that $X$ and $Y$ are connected.
        Notice that for $y\in Y$, $X\times\set y\subseteq X\times Y$ is isomorphic to $X$ via projection.
        This is because the singleton is an affine variety, and is terminal in the category of SWFs; the product of an object with a terminal object is isomorphic
        to the original object via projection.

        Thus $X\times\set y$ is connected, and for the same reason so is $\set x\times Y$ for all $x\in X$.
        If we then fix $y\in Y$ and consider
        $$ Z_x = (X\times\set y) \cup (\set x\times Y), $$
        we see that $Z_x$ is connected as it is the union of non-disjoint connected spaces.
        Then
        $$ X \times Y = \bigcup_x Z_x $$
        is connected since each $Z_x$ is connected and they have non-empty intersection: $\bigcap_xZ_x=X\times\set y$.

        Conversely if $X\times Y$ is connected, then $X=\pi_X(X\times Y)$ is connected as the continuous image of a connected space.
        So is $Y=\pi_Y(X\times Y)$.

        \item The continuous image of an irreducible space is irreducible.
        If $f\colon X\to Y$ is continuous and $X$ is irreducible, then suppose $f(X)=C_1\cup C_2$ is a union of closed subsets.
        Then $X=f^{-1}(C_1\cup C_2)=f^{-1}C_1\cup f^{-1}C_2$ is a proper closed union, so $X=f^{-1}C_i$ for some $i$, which means $f(X)\subseteq C_i$,
        so the original union was not proper; $f(X)$ is irreducible.

        Since $X=\pi_X(X\times Y)$, $X$ is irreducible; the same is true for $Y$.

        \item Suppose $X\times Y=C_1\cup C_2$ be a proper closed union, and let $y\in Y$.
        Then $X\times\set y$ is irreducible as it is isomorphic to $X$; thus it is contained in either $C_1$ or $C_2$.

        Define $F_i=\set{y\in Y}[X\times\set y\subseteq C_i]$, then
        $$ F_i = \bigcap_{x\in X}\pi_Y(C_i\cap(\set x\times Y)) , $$
        since the latter consists of all $y\in Y$ such that for all $x\in X$ there is an $x'\in X$ such that $(x',y)\in C_i\cap(\set x\times Y)$.
        This requires that $x=x'$, so $(x,y)\in C_i$ for all $x\in X$, i.e.\ $X\times\set y\subseteq C_i$.

        Now
        $$ \pi_Y(C_i\cap(\set x\times Y)) = \set{y\in Y}[(x,y)\in C_i] $$
        is the preimage of $C_i$ under $y\mapsto(x,y)$ which is continuous (it is the inverse of $\pi_Y\colon Y\to\set x\times Y$ which we already explained
        is an isomorphism).
        Thus it is closed, so $F_i$ is closed as the intersection of these closed sets.

        Since $Y=F_1\cup F_2$ as for every $y\in Y$, $X\times\set y$ is contained in either $C_1$ or $C_2$.
        Since $Y$ is irreducible, $Y=F_i$ for some $i$.
        Thus $X\times\set y\subseteq C_i$ for all $y\in Y$, thus $X\times Y\subseteq C_i$ as desired.

        \item There is a canonical isomorphism $X^2\times Y^2\cong(X\times Y)^2$ given by permuting coordinates.
        Under this isomorphism, $\Delta(X\times Y)=\Delta(X)\times\Delta(Y)$.
        Thus if $X,Y$ are separated then $\Delta(X)\subseteq X^2$ and $\Delta(Y)\subseteq Y^2$ are closed, and thus so is their product.
        Thus $X\times Y$ is separated.

        Conversely if $X\times Y$ is separated, then for any $y\in Y$, $X\times\set y\subseteq X\times Y$ is a subvariety and thus separated.
        Since $X$ is isomorphic to $X\times\set y$, we see that $X$ is separated.
        Similarly $Y$ is separated.

        \item In lieu of the previous point, we can assume that $X,Y,X\times Y$ are separated.
        If $X$ and $Y$ are complete, then let $W$ be another variety.
        We can factorize the projection $X\times Y\times W\to W$ by
        $$ X\times Y\times W \xtos\proj Y\times W \xtos\proj W . $$
        Since $X,Y$ are complete each projection is a closed map.
        The composition of closed maps is closed, and thus $X\times Y\times W\to W$ is closed, so $X\times Y$ is complete.

        For $y\in Y$, $X\times\set y\subseteq X\times Y$ is a closed subvariety (since $X\subseteq X$ and $\set y\subseteq Y$ are closed, so is their product).
        Thus it is complete (closed subvarieties of complete varieties are complete), and it is isomorphic to $X$ so $X$ is complete.
        Similar for $Y$.
    \eenum
\eenum

\advancecounter{math counter}{1}

\bprob

    Let $U,X,Y$ be algebraic varieties with open embeddings $\alpha\colon U\inj X$ and $\beta\colon U\inj Y$.
    Show that $X\amalg_UY$ is also an algebraic variety.

\eprob

Recall that $X\amalg_UY$ is given as the quotient SWF of $X\amalg Y$ by identifying $\alpha(u)\sim\beta(u)$ for $u\in U$.
Since $\alpha,\beta$ are injective, we only identify an element of $X$ with at most a single element of $Y$ and vice versa.
Indeed, if $x\sim y,y'$ then $x=\alpha(u)$ and $y=\beta(u)$ for some $u\in U$, and also $x=\alpha(u')$ and $y'=\beta(u')$ for another $u'$.
By injectivity, $u=u'$ so $y=y'$.
In particular no two elements of $X$ or two elements of $Y$ are identitified, as this would require them both being identitifed with a single element from
the other set.

Let $\phi_X\colon X\to X\amalg_UY$ be the composition of the inclusion $X\to X\amalg Y$ with the quotient $X\amalg Y\to X\amalg_UY$.
We claim that $\phi_X$ is an open embedding.
It is a morphism of SWFs as the composition of morphisms.
It is injective since as explained, no two elements of $X$ are identified.
Furthermore, $X\to X\amalg Y\to X\amalg_UY$ are both open maps, so their composition $\phi_X$ is as well.
So $\phi_X$ is an injective continuous open map, thus an open embedding.

Now we claim that $\phi_X$ is an isomorphism $X\to\phi_X(X)$.
As an open embedding, it is sufficient to show that $\phi_X^*$ is an isomorphism on all open $V\subseteq\phi_X(X)$.
Let $\pi\colon X\amalg Y\to X\amalg_UY$ be the quotient map, and $\iota_X\colon X\to X\amalg Y$ and $\iota_Y\colon Y\to X\amalg Y$ be the inclusions.
Then
$$ \eqalign{
    \O(V) &= \set{f\colon V\to k}[f\circ\pi\in\O(\pi^{-1}V)]\cr
        &= \set{f\colon V\to k}[f\circ\pi\circ\iota_X\in\O(\iota_X^{-1}\pi^{-1}V),\ f\circ\pi\circ\iota_Y\in\O(\iota_Y^{-1}\pi^{-1}V)]\cr
        &= \set{f\colon V\to k}[\phi_X^*(f)\in\O(\phi_X^{-1}V),\ \phi_Y^*(f)\in\O(\phi_Y^{-1}V)] .
} $$
So now we need to show that if $\phi_X^*(f)\in\O(\phi_X^{-1}V)$ then $\phi_Y^*(f)\in\O(\phi_Y^{-1}V)$.
Now, $\phi_Y^{-1}(\phi_X(X))=\beta(U)$ is isomorphic to $U\cong\alpha(U)$.
Thus $\O(\phi_Y^{-1}V)$ is isomorphic to $\O(\phi_X^{-1}V\cap\alpha(U))$, so we get the desired result.

Thus $\phi_X$ is an isomorphism $X\to\phi_X(X)$.
This means that $X\amalg_UY=\phi_X(X)\cup\phi_Y(Y)$ is an open cover by algebraic varieties.
Thus $X\amalg_UY$ is itself an algebraic variety (the union of affine covers of $\phi_X(X)$ and $\phi_Y(Y)$ is an affine cover of $X\amalg_UY$).

\bprob

    \benum
        \item Let $X$ be a topological space.
        Show that $X$ is Hausdorff iff $\Delta(X)$ is closed in $X\times X$.
        \item Suppose the case in the previous question.
        Show that $X\amalg_UY$ is separated iff $X$ and $Y$ are separated and $U\to X\times Y$ given by $u\mapsto(u,u)$ is a closed embedding.
        \item Find an example of a variety $X$ and two open affine subsets $U,V\subseteq X$ such that $U\cap V$ is not affine.
    \eenum

\eprob

\benum
    \item $\Delta(X)\subseteq X\times X$ is closed iff every $(x,y)\in X\times X-\Delta(X)$ has a neighborhood disjoint from $\Delta(X)$.
    Since product neighborhoods form a basis of $X\times X$, this is iff for all $x\neq y\in X$ there are neighborhoods $x\in U$ and $y\in V$ such that
    $U\times V$ and $\Delta(X)$ are disjoint.
    This is iff $U$ and $V$ are disjoint: a point $(a,b)$ is in the intersection $(U\times V)\cap\Delta(X)$ iff $a\in U,b\in V$ and $a=b$, i.e.\ iff $a=b\in U\cap V$.

    Thus $\Delta(X)\subseteq X\times X$ is closed iff for every $x\neq y\in X$ there are neighborhoods $x\in U$ and $y\in V$ such that $U\cap V=\emptyset$.
    That is, iff $X$ is Hausdorff.

    \item In general the map $u\to(u,u)$ is an embedding since it is continuous, injective (since $U$ embeds into $X$ and $Y$), and open onto its image since it maps $V\subseteq U$ open
    to $V\mapsto (V\times V)\cap\Delta(U)$, where $\Delta(U)=\set{(u,u)}[u\in U]\subseteq X\times Y$.
    And since $U\inj X,Y$ are open embeddings, $V\times V\subseteq X\times Y$ is the product of open subsets and thus open.
    So $u\to(u,u)$ is open onto its image, injective, and continuous; thus an embedding.
    To be a closed embedding we must have that $\Delta(U)\subseteq X\times Y$ is closed.

    Since we can embed $X$ and $Y$ into $Z=X\amalg_UY$ via $\phi_X(X)$ and $\phi_Y(Y)$, and $\phi_X(X)\cap\phi_Y(Y)$ is isomorphic to $U$, this question is equivalent to the following.
    If $Z$ is a variety and $Z=X\cup Y$ is an open cover, then $Z$ is separated iff $X$ and $Y$ are, and $\Delta(X\cap Y)\subseteq X\times Y$ is closed.

    Now notice that
    $$ Z^2 = X^2 \cup (X\times Y) \cup (Y\times X) \cup Y^2 $$
    is an open cover of $Z^2$ (since $X,Y\subseteq Z$ are open).
    Furthermore we have
    $$ \Delta(Z) = \Delta(X) \cup \Delta(Y) \subseteq Z^2 . $$
    A subspace is closed iff it is closed in an open cover, so $\Delta(Z)\subseteq Z^2$ is closed iff its intersections which each of these open subsets is closed (relatively).
    Notice that
    $$ \displaylines{
        \Delta(Z)\cap X^2 = (\Delta(X)\cup\Delta(Y))\cap X^2 = \Delta(X) \cup \Delta(X\cap Y) = \Delta(X) \subseteq X^2,\cr
        \Delta(Z)\cap(X\times Y) = \Delta(X\cap Y) \subseteq X\times Y,\cr
        \Delta(Z)\cap(Y\times X) = \Delta(X\cap Y) \subseteq Y\times X,\cr
        \Delta(Z)\cap Y^2 = \Delta(Y) \subseteq Y^2.
    } $$
    So $\Delta(Z)\subseteq Z^2$ is closed iff all of these sets are closed.
    The first and last lines means that $X$ and $Y$ are separated.
    The second line means that $\Delta(X\cap Y)\subseteq X\times Y$ is closed, so the diagonal morphism is a closed embedding.
    Note that the second and third lines are equivalent: the isomorphism $X\times Y\cong Y\times X$ leaves the diagonal $\Delta(X\cap Y)$ fixed, so
    it is closed in $X\times Y$ iff it is closed in $Y\times X$.

    Thus we get the desired result: $Z$ is separated iff $X$ and $Y$ are, and $\Delta(X\cap Y)\subseteq X\times Y$ is closed.

    \item Consider the diagonal image of $\bA^n$ in $\bA^n\times\bA^n\cong\bA^{2n}$.
    This is $\Delta(\bA^n)=Z(x_1-y_1,\dots,x_n-y_n)$ which is closed in $\bA^{2n}$.
    Furthermore, the diagonal morphism forms an isomorphism $\bA^n\cong\Delta(\bA^n)$.
    Thus if $\Delta(\bA^n-0)\subseteq\bA^{2n}$ were closed, then $\Delta(\bA^n-0)\subseteq\Delta(\bA^n)$ would be closed.
    But this would imply (by taking the inverse of the diagonal morphism, $\Delta(\bA^n)\to\bA^n$) that $\bA^n-0\subseteq\bA^n$ is closed.
    But this is not true, since $\bA^n$ is irreducible and thus connected and $\bA^n-0$ is open (and thus not closed).

    So the diagonal morphism $\bA^n-0\to\bA^{2n}$ is not a closed embedding, so $\bA^n\coprod_{\bA^n-0}\bA^n$ is not separated.

    \item Consider the variety in the previous point.
    Recall that restriction forms an isomorphism $\O(\bA^n)\to\O(\bA^n-0)$, and thus if $\bA^n-0$ were affine then since global sections reflects isomorphisms of affine varieties,
    this would mean inclusion $\bA^n-0\inj\bA^n$ is an isomorphism, which is clearly not the case.

    So if we let $U,V$ be the quotient of the copies of $\bA^n$ in $\bA^n\amalg_{\bA^n-0}\bA^n$, they are affine as they are isomorphic to $\bA^n$ (from the previous question).
    But their intersection is isomorphic to $\bA^n-0$, which is not affine.
\eenum

\advancecounter{math counter}{1}

\bprob

    \benum
        \item Let $\pi\colon\bP^1\times\bA^3\to\bA^3$ be the projection.
        Let the coordinates of $\bP^1$ be $[x:y]$, and the coordinates of $\bA^3$ be $(u,v,w)$.
        Let $Z=Z_{\bP^1\times\bA^3}(x^2+uxy+vy^2,x-wy)$.
        Show explicitly that $\pi(Z)\subseteq\bA^3$ is closed.

        \item Let $X$ be a complete irreducible variety of dimension $\geq1$, and let $x\in X$.
        Show that $X-\set x$ is not universally closed.
    \eenum

\eprob

\benum
    \item Notice that if $x^2+uxy+vy^2=x-wy=0$, then $w=x/y$ since if $y=0$ this would imply $x=0$ a contradiction.
    Thus dividing the first equation by $y^2$ we get
    $$ 0 = \frac{x^2}{y^2} + u\frac xy + v = w^2 + uw + v . $$
    And so we see that
    $$ \pi(Z) = Z_{\bA^3}(w^2+uw+v) . $$
    This is because if $(u,v,w)\in\pi(Z)$ then there is some $[x:y]$ such that $([x:y],(u,v,w))\in Z$ which implies $w^2+uw+v=0$ as we showed.
    Conversely if $w^2+uw+v=0$ then choosing $y=1$ and $x=w$, we get that $([x:y],(u,v,w))\in Z$ and so $(u,v,w)\in\pi(Z)$.

    So $\pi(Z)$ is indeed closed, as desired.

    \item Consider the projection $\pi\colon(X-\set x)\times X\to X$.
    Since $X$ is complete and thus separated, $\Delta(X)\subseteq X\times X$ is closed.
    Thus $\Delta(X)\cap(X-\set x)\times X\subseteq(X-\set x)\times X$ is closed, and it is equal to
    $$ \Delta(X)\cap(X-\set x)\times X = \set{(y,y)\in Y\times Y}[y\neq x] = \Delta(X-\set x) . $$
    But $\pi(\Delta(X-\set x))=X-\set x\subseteq X$ is not closed, as it is open in $X$ and $X$ is irreducible and thus connected.
\eenum

\bprob[title={The rigidity lemma}]

    Let $X,Y,Z$ separated irreducible varieties such that $X$ is further complete.
    Let $f\colon X\times Y\to Z$ be a morphism such that there are $x_0\in X,y_0\in Y,z_0\in Z$ where
    $$ f(x_0\times Y) = z_0 = f(X\times y_0) . $$
    Show that $f$ is the constant function $z_0$.

\eprob

Define
$$ S = \set{y\in Y}[f(X\times y)=z_0] , $$
which is non-empty as it contains $y_0$.
Furthermore we can write it as
$$ S = \bigcap_{x\in X}\set{y\in Y}[f(x,y)=z_0] = \bigcap_{x\in X}f_x^{-1}(z_0) $$
where $f_x(y)=f(x,y)$ is a morphism $Y\to Z$, so that $f_x^{-1}z_0$ is closed for all $x\in X$ and thus $S$ is closed.

Now consider the graph of $f$, $\Gamma(f)\subseteq X\times Y\times Z$.
Since $Z$ is separated, $\Gamma(f)$ is closed, and since $X$ is complete, $W=\pi(\Gamma(f))\subseteq Y\times Z$ is closed.
Consider also the projection $\pi_Y\colon W\to Y$, which has fibers
$$ \pi_Y^{-1}y = \set{(y,z)\in Y\times Z}[f(x,y)=z\hbox{ for some $x\in X$}] = \set y\times f(X\times y) . $$
Notice that $f(X\times\set y)$ is a point iff $\pi_Y^{-1}y$ is a point.
Since $X$ is irreducible, so is $X\times y$, and thus so is $f(X\times y)$ (the continuous image of an irreducible space is irreducible).
Thus $\pi_Y^{-1}y$ is irreducible, and so it is a point iff it has dimension zero.
Furthermore as it is irreducible, for any $(y,z)\in\pi_Y^{-1}y$, $\dim_{(y,z)}\pi_Y^{-1}y=\dim\pi_Y^{-1}y$.

Now $f(X\times y)=z_0$ iff $f(X\times y)$ is a point, as $X\times y$ contains $(x_0,y)$ whose image is $z_0$.
Thus
$$ S = \set{y\in Y}[f(X\times y)=z_0] = \set{y\in Y}[\hbox{$f(X\times y)$ is a point}] = \set{y\in Y}[\dim\pi_Y^{-1}y=0] . $$
Now we know that
$$ \set{(y,z)\in W}[\dim_{(y,z)}\pi_Y^{-1}y\geq1] = \set{(y,z)\in W}[\dim\pi_Y^{-1}y\geq1]\subseteq W $$
is closed.
Thus $\set{(y,z)\in W}[\dim\pi_Y^{-1}y=0]=S\times Z\subseteq W$ is open, and thus $S$ is open.

Since $S\subseteq Y$ is nonempty open and closed, and $Y$ is irreducible, $S$ must be all of $Y$.
So $f(X\times y)=z_0$ for all $y\in Y$, i.e.\ $f(x,y)=z_0$ for all $x\in X,y\in Y$ as desired.

\bprob

    An {\emphcolor algebraic group} is a group object $G$ in the category of varieties.
    That is, it is a variety $G$ with morphisms $m\colon G\times G\to G$ (multiplication) and $i\colon G\to G$ (inversion) satisfying the group axioms.
    $G$ is an {\emphcolor Abelian variety} if it is complete and irreducible.
    Show that an Abelian variety is an Abelian group.

\eprob

Consider the commutator morphism $C\colon G\times G\to G$ given by $C(g,h)=ghg^{-1}h^{-1}$.
$G$ is an Abelian iff $C$ is trivial, i.e.\ $C(g,h)=1$ for all $g,h\in G$.
Notice that $C$ is a morphism of varieties, as it can be written as the composition
$$ G\times G \xtos{(m,i\circ m)} G\times G \xtos m G . $$
The first arrow is a morphism as each coordinate is (and by the universal property of products).

Now we have for all $g\in G$, $C(1,g)=C(g,1)=1$.
Thus by the rigidity lemma $C$ is the constant map $1$, meaning $G$ is indeed Abelian.

\bye

